%%============================================================================
%% POLÍTICA DE PRIVACIDADE — Grupo de Usuários Linux e Software Livre
%%                              Estatística e Ciência de Dados - UFPR
%%============================================================================
%% Documento institucional — versão formal
%% Compilar com: pdflatex politica-privacidade.tex  (ou xelatex / lualatex)
%%============================================================================

\documentclass[
  12pt,          % tamanho base da fonte
  a4paper,       % formato de papel
  oneside,       % impressão frente única
  brazil         % idioma principal: português do Brasil
]{article}

%---------------------------------------------------------------------------
% Pacotes essenciais
%---------------------------------------------------------------------------
\usepackage[T1]{fontenc}            % codificação de fonte
\usepackage[utf8]{inputenc}         % entrada UTF-8 (pdflatex)
\usepackage{lmodern}                % fontes Latin Modern (vetoriais)
\usepackage{microtype}              % microtipografia (melhora legibilidade)

\usepackage[brazil]{babel}          % hifenização e nomes de seções em PT-BR

\usepackage[
  margin=2.5cm,
  top=3cm,
  bottom=3cm,
  headheight=25pt
]{geometry}                         % margens do documento

\usepackage{setspace}               % controle de espaçamento entre linhas
\usepackage{xcolor}                 % suporte a cores
\usepackage{titlesec}               % formatação de títulos de seções
\usepackage{titling}                % formatação do título do documento
\usepackage{fancyhdr}               % cabeçalhos e rodapés personalizados
\usepackage{enumitem}               % personalização de listas
\usepackage{lettrine}               % capitulares (letra inicial decorada)
\usepackage{graphicx}               % suporte a imagens (se necessário)
\usepackage{hyperref}               % hiperlinks e metadados do PDF
\usepackage{bookmark}               % bookmarks do PDF mais limpos

%---------------------------------------------------------------------------
% Cores institucionais
%---------------------------------------------------------------------------
\definecolor{gulecdBlue}{HTML}{1B3A5C}   % azul escuro institucional
\definecolor{gulecdAccent}{HTML}{2E86AB} % azul de destaque
\definecolor{gulecdGray}{HTML}{555555}   % cinza para texto secundário
\definecolor{gulecdLight}{HTML}{F5F7FA}  % fundo claro para caixas
\definecolor{gulecdRule}{HTML}{D0D5DD}   % cor de filetes e réguas

%---------------------------------------------------------------------------
% Metadados do PDF
%---------------------------------------------------------------------------
\hypersetup{
  pdftitle    = {Política de Privacidade — GULECD/UFPR},
  pdfauthor   = {Grupo de Usuários Linux - Estatística e Ciência de Dados - UFPR},
  pdfsubject  = {Política de Privacidade institucional do GULECD},
  pdfkeywords = {política de privacidade, proteção de dados, LGPD, software livre, código aberto, ciência de dados, UFPR},
  colorlinks  = true,
  linkcolor   = gulecdAccent,
  urlcolor    = gulecdAccent,
  citecolor   = gulecdAccent
}

%---------------------------------------------------------------------------
% Configurações tipográficas
%---------------------------------------------------------------------------
\onehalfspacing                     % espaçamento entre linhas = 1,5
\setlength{\parindent}{1.5cm}       % recuo de parágrafo
\setlength{\parskip}{0.5em}         % espaço entre parágrafos

%---------------------------------------------------------------------------
% Formato do título do documento
%---------------------------------------------------------------------------
\pretitle{%
  \begin{center}
  \vspace*{1cm}
  \rule{\textwidth}{2pt}\vspace{0.3cm}\\
  \LARGE\bfseries\color{gulecdBlue}
}
\posttitle{%
  \\\vspace{0.3cm}
  \rule{\textwidth}{2pt}
  \end{center}
  \vspace{0.5cm}
}
\preauthor{\begin{center}\large\color{gulecdGray}}
\postauthor{\end{center}}
\predate{\begin{center}\small\color{gulecdGray}}
\postdate{\end{center}}

%---------------------------------------------------------------------------
% Formato das seções
%---------------------------------------------------------------------------
\titleformat{\section}
  [block]
  {\normalfont\Large\bfseries\color{gulecdBlue}}
  {\thesection.}{0.8em}{}
  [\vspace{0.2em}{\color{gulecdRule}\rule{\linewidth}{0.6pt}}]

\titleformat{\subsection}
  [block]
  {\normalfont\large\bfseries\color{gulecdAccent}}
  {\thesubsection.}{0.8em}{}

\titlespacing*{\section}
  {0pt}{2.5ex plus 1ex minus .2ex}{1.5ex plus .2ex}
\titlespacing*{\subsection}
  {0pt}{2ex plus 1ex minus .2ex}{1ex plus .2ex}

%---------------------------------------------------------------------------
% Cabeçalho e rodapé
%---------------------------------------------------------------------------
\pagestyle{fancy}
\fancyhf{}
\fancyhead[L]{\small\color{gulecdGray}\textsc{GULECD/UFPR}}
\fancyhead[R]{\small\color{gulecdGray}\textit{Política de Privacidade}}
\fancyfoot[C]{\small\color{gulecdGray}---\ \thepage\ ---}
\renewcommand{\headrulewidth}{0.4pt}
\renewcommand{\headrule}{%
  \hbox to\headwidth{%
    \color{gulecdRule}\leaders\hrule height \headrulewidth\hfill}}

% Remove cabeçalho/rodapé da primeira página (será redefinido manualmente)
\fancypagestyle{plain}{%
  \fancyhf{}%
  \fancyfoot[C]{\small\color{gulecdGray}---\ \thepage\ ---}%
  \renewcommand{\headrulewidth}{0pt}%
}

%---------------------------------------------------------------------------
% Ambiente personalizado para listas
%---------------------------------------------------------------------------
\newlist{privacidade}{itemize}{1}
\setlist[privacidade]{
  label      = {\color{gulecdAccent}\textbullet},
  leftmargin = 1.5em,
  itemsep    = 0.3em,
  topsep     = 0.4em
}

%---------------------------------------------------------------------------
% Informações do documento
%---------------------------------------------------------------------------
\title{Política de Privacidade\\[0.3em]
  \Large\mdseries\color{gulecdGray}%
  Grupo de Usuários Linux e \textit{Software} Livre\\Estatística e Ciência de Dados - UFPR}
\date{\today}

%============================================================================
% INÍCIO DO DOCUMENTO
%============================================================================
\begin{document}

%--- Capa / Título principal -------------------------------------------------
\thispagestyle{empty}
\maketitle

\newpage

%--- Seção 1: Introdução ----------------------------------------------------
\section{Introdução}

\lettrine[lines=3, lhang=0.1, nindent=0em, slope=-0.5em]{E}{sta} Política de Privacidade estabelece as diretrizes para coleta, uso, armazenamento
 e proteção de dados pessoais no âmbito do GULECD/UFPR.
O grupo compromete-se a adotar práticas alinhadas aos princípios de privacidade, transparência e minimização de dados, em conformidade com a legislação brasileira aplicável,
em especial a Lei Geral de Proteção de Dados (LGPD – Lei nº 13.709/2018).

%--- Seção 2: Princípios Gerais ---------------------------------------------
\section{Princípios Gerais}

O GULECD/UFPR adota os seguintes princípios no tratamento de dados:

\begin{privacidade}
	\item \textbf{Finalidade específica}: dados são coletados apenas para objetivos claramente definidos
	\item \textbf{Minimização}: coleta-se apenas o mínimo necessário
	\item \textbf{Transparência}: os participantes devem saber quais dados são coletados e por quê
	\item \textbf{Segurança}: dados são protegidos com boas práticas técnicas e organizacionais
	\item \textbf{Não compartilhamento}: dados não são repassados a terceiros, exceto quando exigido por lei
\end{privacidade}

%--- Seção 3: Coleta de Dados -----------------------------------------------
\section{Coleta de Dados}

Os dados coletados podem incluir:

\begin{privacidade}
	\item nome completo
	\item endereço de email
	\item telefone
	\item vínculo institucional (quando aplicável)
	\item matrícula (quando necessário para emissão de certificados ou reconhecimento de participação pela UFPR)
	\item outras informações estritamente necessárias ao contexto da atividade
\end{privacidade}

Nenhum dado será coletado sem um \textbf{objetivo específico, explícito e legítimo}.

%--- Seção 4: Uso dos Dados -------------------------------------------------
\section{Uso dos Dados}

Os dados coletados serão utilizados exclusivamente para fins internos do grupo, incluindo:

\begin{privacidade}
	\item organização e gestão de atividades
	\item comunicação com membros e participantes
	\item registro histórico e institucional
	\item operação de projetos e iniciativas
	\item inscrição em atividades ou eventos
	\item adesão a documentos institucionais
	\item participação em projetos
\end{privacidade}

Os dados \textbf{não serão utilizados para fins comerciais} ou qualquer finalidade externa ao escopo do grupo.

%--- Seção 5: Compartilhamento de Dados -------------------------------------
\section{Compartilhamento de Dados}

O GULECD/UFPR não compartilha dados pessoais com terceiros.

Exceções ocorrem apenas quando:

\begin{privacidade}
	\item houver \textbf{obrigação legal}, conforme legislação brasileira
	\item houver \textbf{requisição formal de autoridade competente}
\end{privacidade}

Fora dessas situações, os dados permanecem sob uso interno exclusivo do grupo.

%--- Seção 6: Armazenamento e Segurança -------------------------------------
\section{Armazenamento e Segurança}

Os dados serão armazenados utilizando as melhores práticas disponíveis, incluindo:

\begin{privacidade}
	\item uso de plataformas confiáveis
	\item controle de acesso restrito
	\item armazenamento seguro (criptografia quando aplicável)
	\item revisão periódica de acessos
\end{privacidade}

O grupo compromete-se a reduzir riscos de:

\begin{privacidade}
	\item acesso não autorizado
	\item vazamento de dados
	\item uso indevido
\end{privacidade}

%--- Seção 7: Responsabilidade ----------------------------------------------
\section{Responsabilidade}

A responsabilidade pelo cumprimento desta política recai sobre a \textbf{Coordenação do GULECD/UFPR}, que deve:

\begin{privacidade}
	\item garantir a aplicação das diretrizes aqui estabelecidas
	\item supervisionar o uso adequado dos dados
	\item responder a solicitações relacionadas à privacidade
\end{privacidade}

%--- Seção 8: Direitos dos Titulares ----------------------------------------
\section{Direitos dos Titulares}

Qualquer membro ou participante poderá, a qualquer momento:

\begin{privacidade}
	\item solicitar acesso aos seus dados
	\item solicitar correção de dados incorretos
	\item solicitar \textbf{remoção completa de seus dados}
\end{privacidade}

As solicitações devem ser atendidas em prazo razoável, respeitando limitações operacionais e obrigações legais.

%--- Seção 9: Retenção de Dados ---------------------------------------------
\section{Retenção de Dados}

Os dados serão mantidos apenas pelo tempo necessário para cumprir suas finalidades.

Quando não forem mais necessários:

\begin{privacidade}
	\item serão removidos, ou
	\item anonimizados, quando aplicável
\end{privacidade}

%--- Seção 10: Atualizações desta Política ----------------------------------
\section{Atualizações deste Documento}

Este documento deve ser entendido como um instrumento vivo, passível de aprimoramento conforme a evolução do grupo, mudanças operacionais ou alterações legais. 
Sua aplicação deve sempre priorizar os princípios fundamentais de abertura, participação, eficiência e responsabilidade coletiva.

Novas versões deste documento deverão ser arquivadas juntamente com as versões anteriores no repositório \textit{documentos-referendados} e numeradas em acordo
 com o esquema padrão de \textit{releases} como XX.yy, onde XX refere-se a edições substanciais e yy à numeração em versões de ajuste. 
 
Também deverá ser mantida ativa a versão em edição no repositório \textit{documents-dev} a fim de prover o rastreamento de todas as mudanças entre as versões do documento,

Assim, determina-se que:
\begin{conduta}
	\item Este regulamento poderá ser revisado mediante discussão coletiva
	\item Situações não previstas serão resolvidas pela Coordenação vigente
	\item As atualizações serão comunicadas de forma transparente aos membros
\end{conduta}

%--- Seção 11: Disposições Finais -------------------------------------------
\section{Disposições Finais}

O GULECD/UFPR reafirma seu compromisso com a privacidade e proteção de dados, reconhecendo que a confiança dos membros e participantes é fundamental para o funcionamento 
e crescimento do grupo.


%============================================================================
\end{document}
%============================================================================
